\documentclass[a4paper,12pt]{article}

\usepackage{fontspec}
\setmainfont{Times New Roman}
\usepackage[russian]{babel}
\usepackage{graphicx, float}
\usepackage{geometry}
\geometry{left=3cm,right=1.5cm,top=2cm,bottom=2cm}
\usepackage{setspace}
\onehalfspacing
\setlength{\parindent}{1.25cm}

\begin{document}

\begin{titlepage}
  \thispagestyle{empty}
  \begin{center}

    
\includegraphics[width=0.18\textwidth]{../assets/gerb.png}

    \vspace{0.5cm}

    {\small
      МИНИСТЕРСТВО НАУКИ И ВЫСШЕГО ОБРАЗОВАНИЯ\\
      РОССИЙСКОЙ ФЕДЕРАЦИИ
    }

    \vspace{0.3cm}

    {\small
      Федеральное государственное бюджетное образовательное учреждение\\
      высшего образования
    }

    \vspace{0.3cm}

    {\large \textbf{«МИРЭА --- Российский технологический университет»}}

    {\normalsize РТУ МИРЭА}

    {\normalsize Институт искусственного интеллекта (ИИИ)}\\
    {\normalsize Кафедра промышленной информатики (ПИ)}

    \vspace{2cm}

    {\LARGE \textbf{ОТЧЁТ}}

    \vspace{0.3cm}

    {\Large \textbf{ПО ПРАКТИЧЕСКОЙ РАБОТЕ № 24}}

    \vspace{0.4cm}

    {\large по дисциплине «Информатика»}

    \vspace{2.5cm}

    \begin{tabular}{ll}
      Выполнил студент группы КВБО-11-25 & Шапаренко Ф.А. \\
      Принял ассистент: & Беднов Г.А. \\
    \end{tabular}

    \vfill

    Практическую работу выполнил: «04» мая 2026 г.\\
    «Зачтено» \rule{1cm}{0.4pt} 2026 г.

    \vspace{1cm}

    Москва 2026

  \end{center}
\end{titlepage}

\noindent\textbf{Тема:} «Описательная статистика»\\
\noindent\textbf{Цель работы:} Научится вычислять параметры,
определяемые инструментом  Описательная статистика.

\begin{figure}[h]
  \centering
  \includegraphics[width=0.5\textwidth]{/home/shafed/projects/02-computer-science/defend-4/assets/24.png}
\end{figure}

Среднее значение (51,85): среднее арифметическое всех 100 чисел. Оно
показывает общий “центр” выборки.

Стандартная ошибка (2,78): показатель того, насколько точно среднее
значение выборки (51,85) отражает реальное среднее значение всей совокупности.

Медиана (54): центральное значение в упорядоченном наборе данных.
Половина чисел меньше или равна 54, а половина больше или равна 54.

Мода (19): число, которое встречается в наборе чаще всего.

Первый квартиль (30,75): значение, ниже которого находится примерно
25% всех данных.

Третий квартиль (77): значение, ниже которого находится примерно 75%
% всех данных.

Стандартное отклонение (27,85) и Дисперсия (775,60): эти цифры
показывают, насколько сильно числа «разбросаны» относительно среднего
значения. Чем они выше, тем больше разброс.

Асимметричность (-0,04): значение почти равно нулю, значит данные
распределены почти симметрично, без заметного перекоса вправо или влево.

Эксцесс (-1,09): отрицательное значение указывает на то, что
распределение чисел более «плоское», чем нормальное, то есть нет ярко
выраженного пика в центре.

Размах (98): разница между максимальным и минимальным значением в выборке.

Минимум (2): самое маленькое число в наборе данных.

Максимум (100): самое большое число в наборе данных.

Сумма (5185): сумма всех 100 значений выборки.

Количество (100): общее число значений в выборке.

Наименьший(2) (2): второе наименьшее значение в наборе данных. Оно
тоже равно 2, потому что минимальное значение встречается больше одного раза.

Наибольший(2) (99): второе наибольшее значение в наборе данных после
максимума 100.

\end{document}
